\documentclass[11pt]{article}

\usepackage[margin=1in]{geometry}
\usepackage{amsmath}
\usepackage{amssymb}
\usepackage{booktabs}
\usepackage{hyperref}
\usepackage{xcolor}
\usepackage{listings}

\hypersetup{
  colorlinks = true,
  linkcolor = blue,
  urlcolor = blue
}

\lstdefinestyle{Rstyle}{
  language = R,
  basicstyle = \ttfamily\small,
  keywordstyle = \color{blue},
  commentstyle = \color{gray},
  stringstyle = \color{teal},
  showstringspaces = false,
  columns = fullflexible,
  breaklines = true,
  frame = single
}
\lstset{style = Rstyle}

\title{\textbf{MALC Package Manual}\\
\large Density Estimation from Ordered Log-Concave Discrete Data}
\author{Furkan Danisman}
\date{Version 0.1.0}

\begin{document}
\maketitle

\tableofcontents
\newpage

\section{Overview}

\texttt{MALC} implements density estimation from observed frequencies on a
known grid. The package is designed for application settings where the raw
sample is already aggregated into bin counts. The main estimator is
\texttt{MALC()}, and the package also includes the three comparison methods
used in the simulation study:

\begin{enumerate}
  \item \texttt{BK2002()}
  \item \texttt{BinnedNP()}
  \item \texttt{KernSmooth()}
\end{enumerate}

All fitting functions use the same basic input structure:

\begin{itemize}
  \item \texttt{counts}: observed frequencies for each grid interval.
  \item \texttt{grid}: numeric grid breaks with length \texttt{length(counts) + 1}.
\end{itemize}

The MALC fit supports density, distribution, quantile, mean, and plotting
methods. Kernel comparison fits support L2 evaluation through \texttt{Eval()}
and visualization through \texttt{Plot\_Kernel()}.

\section{Installation}

The package can be installed from GitHub while the package branch is under
development:

\begin{lstlisting}
install.packages("remotes")
remotes::install_github("FurkanDanisman/MALC")
\end{lstlisting}

Then load the package:

\begin{lstlisting}
library(MALC)
\end{lstlisting}

\section{Input Format}

\texttt{MALC()} does not accept raw observations. It accepts only observed
frequencies and grid breaks.

\begin{lstlisting}
set.seed(1)
x <- rnorm(1000)
grid <- seq(-4, 4, by = 0.5)
counts <- as.integer(table(cut(x, breaks = grid, include.lowest = TRUE)))
\end{lstlisting}

In direct package usage, the preprocessing step above can be done by the user
or by an application before calling the package.

\section{MALC Estimator}

\subsection{Fit}

\begin{lstlisting}
fit <- MALC(counts, grid)
fit_s <- MALC(counts, grid, smooth = TRUE)
\end{lstlisting}

\subsection*{Usage}

\begin{lstlisting}
MALC(counts, grid, smooth = FALSE, B = 10000, alpha = 2,
         seed = 20180621, log_conc = TRUE)
\end{lstlisting}

\subsection*{Arguments}

\begin{tabular}{ll}
\toprule
Argument & Description \\
\midrule
\texttt{counts} & Integer observed frequencies for each bin. \\
\texttt{grid} & Numeric bin breaks. Must have length \texttt{length(counts) + 1}. \\
\texttt{smooth} & Logical. If \texttt{TRUE}, downstream density functions use the smoothed fit. \\
\texttt{B} & Number of Monte Carlo draws used internally. \\
\texttt{alpha} & Shape parameter used in the beta perturbation. \\
\texttt{seed} & Random seed used internally. \\
\texttt{log\_conc} & Passed to \texttt{logcondens::logConDens()}. \\
\bottomrule
\end{tabular}

\section{MALC Methods}

\subsection{\texttt{dmalc()}}

Evaluate the fitted MALC density.

\begin{lstlisting}
x_eval <- seq(min(grid), max(grid), length.out = 200)
d_hat <- dmalc(fit, x_eval)
\end{lstlisting}

\subsection*{Usage}

\begin{lstlisting}
dmalc(object, x)
\end{lstlisting}

\subsection{\texttt{pmalc()}}

Evaluate the fitted distribution function numerically.

\begin{lstlisting}
pmalc(fit, c(-1, 0, 1))
\end{lstlisting}

\subsection*{Usage}

\begin{lstlisting}
pmalc(object, q, length_grid = 1001)
\end{lstlisting}

\subsection{\texttt{qmalc()}}

Evaluate fitted quantiles numerically.

\begin{lstlisting}
qmalc(fit, c(0.25, 0.50, 0.75))
\end{lstlisting}

\subsection*{Usage}

\begin{lstlisting}
qmalc(object, p, length_grid = 1001)
\end{lstlisting}

\subsection{\texttt{mean\_malc()}}

Return the EM-step mean estimate stored in the fitted object. This is not a
smoothed/non-smoothed choice; it is the mean estimate from the MALC fitting
step.

\begin{lstlisting}
mean_malc(fit)
\end{lstlisting}

\subsection*{Usage}

\begin{lstlisting}
mean_malc(object)
\end{lstlisting}

\subsection{\texttt{Plot()}}

Plot a fitted MALC density. The plotted density follows the \texttt{smooth}
choice stored in the fitted object.

\begin{lstlisting}
Plot(fit)
Plot(fit_s, col = "red")
\end{lstlisting}

\subsection*{Usage}

\begin{lstlisting}
Plot(object, ...)
plot(x, eval_grid = NULL, add = FALSE, xlab = "x",
     ylab = "Density", col = "navy", lwd = 2, ...)
\end{lstlisting}

\section{Comparison Methods}

\subsection{\texttt{BK2002()}}

Fit the BK2002 comparison estimator from observed frequencies and grid breaks.

\begin{lstlisting}
bk <- BK2002(counts, grid)
\end{lstlisting}

\subsection*{Usage}

\begin{lstlisting}
BK2002(counts, grid, h = NULL, bw_B = 50, bw_seed = 1,
       bw_method = c("ste", "dpi"), max_iter = 200,
       tol_L1 = 1e-6, verbose = FALSE)
\end{lstlisting}

\subsection{\texttt{BinnedNP()}}

Fit the BinnedNP comparison estimator. The bandwidth choice is set at fitting
time and stored in the fitted object.

\begin{lstlisting}
bn <- BinnedNP(counts, grid, bandwidth = "boot")
bn_plugin <- BinnedNP(counts, grid, bandwidth = "plugin")
\end{lstlisting}

\subsection*{Usage}

\begin{lstlisting}
BinnedNP(counts, grid, bandwidth = c("boot", "plugin"),
         seed_start = 1L, max_tries = 200L)
\end{lstlisting}

\subsection{\texttt{KernSmooth()}}

Fit the KernSmooth comparison estimator from observed frequencies and grid
breaks.

\begin{lstlisting}
ks <- KernSmooth(counts, grid)
\end{lstlisting}

\subsection*{Usage}

\begin{lstlisting}
KernSmooth(counts, grid, seed = 1, scalest = "minim",
           level = 2L, kernel = "normal", canonical = FALSE,
           truncate = TRUE)
\end{lstlisting}

\section{Kernel Plotting}

\texttt{Plot\_Kernel()} plots any fitted comparison-method object returned by
\texttt{BK2002()}, \texttt{BinnedNP()}, or \texttt{KernSmooth()}.

\begin{lstlisting}
Plot_Kernel(bk)
Plot_Kernel(bn, add = TRUE, col = "red")
Plot_Kernel(ks, add = TRUE, col = "blue")
\end{lstlisting}

\subsection*{Usage}

\begin{lstlisting}
Plot_Kernel(object, eval_grid = NULL, add = FALSE,
            xlab = "x", ylab = "Density", col = "darkgreen",
            lwd = 2, ...)
\end{lstlisting}

\section{L2 Evaluation}

\texttt{Eval()} computes the numerical L2 distance between a fitted density
estimator and a supplied true density function. It accepts both MALC fits
and kernel-method fits.

\begin{lstlisting}
true_pdf <- function(z) dnorm(z)

Eval(fit, true_pdf)
Eval(fit_s, true_pdf)
Eval(bk, true_pdf)
Eval(bn, true_pdf)
Eval(ks, true_pdf)
\end{lstlisting}

\subsection*{Usage}

\begin{lstlisting}
Eval(object, pdf, length_grid = 1001)
\end{lstlisting}

\subsection*{Arguments}

\begin{tabular}{ll}
\toprule
Argument & Description \\
\midrule
\texttt{object} & Fitted object from \texttt{MALC()}, \texttt{BK2002()}, \texttt{BinnedNP()}, or \texttt{KernSmooth()}. \\
\texttt{pdf} & True density function used for comparison. \\
\texttt{length\_grid} & Number of numerical evaluation points. \\
\bottomrule
\end{tabular}

\section{Complete Example}

\begin{lstlisting}
library(MALC)

set.seed(1)
x <- rnorm(1000)
grid <- seq(-4, 4, by = 0.5)
counts <- as.integer(table(cut(x, breaks = grid, include.lowest = TRUE)))

fit <- MALC(counts, grid)
fit_s <- MALC(counts, grid, smooth = TRUE)

dmalc(fit, seq(-2, 2, length.out = 5))
pmalc(fit, c(-1, 0, 1))
qmalc(fit, c(0.25, 0.50, 0.75))
mean_malc(fit)
Plot(fit)

bk <- BK2002(counts, grid)
bn <- BinnedNP(counts, grid, bandwidth = "boot")
ks <- KernSmooth(counts, grid)

Plot_Kernel(bk)
Plot_Kernel(bn, add = TRUE, col = "red")
Plot_Kernel(ks, add = TRUE, col = "blue")

true_pdf <- function(z) dnorm(z)
Eval(fit, true_pdf)
Eval(fit_s, true_pdf)
Eval(bk, true_pdf)
Eval(bn, true_pdf)
Eval(ks, true_pdf)
\end{lstlisting}

\section{Exported Functions}

\begin{itemize}
  \item \texttt{MALC()}
  \item \texttt{dmalc()}
  \item \texttt{pmalc()}
  \item \texttt{qmalc()}
  \item \texttt{mean\_malc()}
  \item \texttt{Plot()}
  \item \texttt{BK2002()}
  \item \texttt{BinnedNP()}
  \item \texttt{KernSmooth()}
  \item \texttt{Plot\_Kernel()}
  \item \texttt{Eval()}
\end{itemize}

\end{document}
